% Espectro instantáneo frente a cascada (waterfall): dos vistas del mismo flujo I/Q.
% Ilustración esquemática; no es una captura de SDRangel.
\begin{tikzpicture}[x=1mm,y=1mm,font=\sffamily\scriptsize]
  % espectro
  \draw[MAMuted] (0,30) rectangle (70,52);
  \draw[MAInk,line width=0.7pt] plot[smooth] coordinates {(0,33) (8,34) (12,33.5) (16,40) (19,46) (22,40) (26,33.6) (34,34) (40,33) (45,37) (47,39) (49,37) (54,33.5) (62,34) (70,33)};
  \node[anchor=west,font=\sffamily\scriptsize\bfseries] at (0,55) {Espectro: potencia frente a frecuencia (un instante)};
  \node[rotulo,anchor=west] at (72,46) {eje vertical: dB};
  \node[rotulo,anchor=west] at (72,40) {una curva por FFT};
  % waterfall
  \foreach \y in {0,2,...,26} {
    \fill[MAPrimarySoft] (0,\y) rectangle (70,\y+2);
    \fill[MAPrimaryDark!70] (16,\y) rectangle (22,\y+2);
  }
  \foreach \y in {4,6,14,16,18,24} \fill[MASecondary!60] (45,\y) rectangle (49,\y+2);
  \draw[MAMuted] (0,0) rectangle (70,28);
  \node[anchor=west,font=\sffamily\scriptsize\bfseries] at (0,-4) {Cascada: cada fila es un espectro; el color es la potencia};
  \draw[flecha] (-3,28) -- node[left,rotulo,align=right]{tiempo\\(las filas\\bajan)} (-3,2);
  \node[rotulo,anchor=west,align=left] at (72,22) {traza continua $\Rightarrow$ emisión permanente\\(p.\ ej.\ radiodifusión FM)};
  \node[rotulo,anchor=west,align=left] at (72,11) {trazos cortos $\Rightarrow$ emisión intermitente\\(ráfagas, como 1090 MHz)};
  \draw[flechadash] (71.5,22) -- (23,20);
  \draw[flechadash] (71.5,11) -- (50,15);
\end{tikzpicture}
